\section{README to LKMM's documentation directory}

\begin{Note}
  Pretty-printed \path{tools/memory-model/Documentation/README}.
\end{Note}

It has been said that successful communication requires first identifying
what your audience knows and then building a bridge from their current
knowledge to what they need to know.
Unfortunately, the expected Linux-kernel memory model (LKMM) audience
might be anywhere from noviceto expert both in kernel hacking and in
understanding LKMM.

This document therefore points out a number of places to start reading,
depending on what you know and what you would like to learn.
Please note that the documents later in this list assume that the reader
understands the material provided by documents earlier in this list.

If LKMM-specific terms lost you,
\hyperref[sec:docs:Glossaries to LKMM]{\path{glossary.txt}}
might help you.

\begin{itemize}
  \item You are new to Linux-kernel concurrency:
	\hyperref[sec:docs:Quick and simple guide to LKMM]{\path{simple.txt}}

  \item You have some background in Linux-kernel concurrency, and would
	like an overview of the types of low-level concurrency primitives
	that the Linux kernel provides:
	\hyperref[sec:docs:Overview of memory-ordering operations]{
	  \path{ordering.txt}}

	Here, ``low level'' means atomic operations to single variables.

  \item You are familiar with the Linux-kernel concurrency primitives
	that you need, and just want to get started with LKMM litmus
	tests:
	\hyperref[sec:docs:Linux-kernel memory model litmus tests]{
	  \path{litmus-tests.txt}}

  \item You need to locklessly access shared variables that are otherwise
	protected by a lock:
	\hyperref[sec:docs:Locking]{\path{locking.txt}}

	This \path{locking.txt} file expands on
	\cref{sec:docs:recipes:Locking}
	in \path{recipes.txt}, but it is self-contained.

  \item You are familiar with Linux-kernel concurrency, and would
	like a detailed intuitive understanding of LKMM, including
	situations involving more than two threads:
	\hyperref[sec:docs:Recipes]{\path{recipes.txt}}

  \item You would like a detailed understanding of what your compiler can
	and cannot do to control dependencies:
	\hyperref[sec:docs:Control dependencies]{
	  \path{control-dependencies.txt}}

  \item You would like to mark concurrent normal accesses to shared
	variables so that intentional ``racy'' accesses can be properly
	documented, especially when you are responding to complaints
	from KCSAN\@:
	\hyperref[sec:docs:Marking shared-memory accesses]{
	  \path{access-marking.txt}}

  \item You are familiar with Linux-kernel concurrency and the use of
	LKMM, and would like a quick reference:
	\hyperref[sec:docs:Quick-reference guide to the Linux-kernel memory model]{
	  \path{cheatsheet.txt}}

  \item You are familiar with Linux-kernel concurrency and the use
	of LKMM, and would like to learn about LKMM's requirements,
	rationale, and implementation:
	\hyperref[sec:docs:Explanation of the Linux-Kernel Memory Consistency Model]{
	  \path{explanation.txt}}
	and \hyperref[sec:docs:Herd representation]{
	  \path{herd-representation.txt}}

  \item You are interested in the publications related to LKMM, including
	hardware manuals, academic literature, standards-committee
	working papers, and LWN articles:
	\hyperref[sec:docs:References]{
	  \path{references.txt}}
\end{itemize}


\subsection{Description of files}

\begin{description}[style=nextline]
  \item[\path{README}]
	This file.

  \item[{\hyperref[sec:docs:Marking shared-memory accesses]{
	  \path{access-marking.txt}}}]
	Guidelines for marking intentionally concurrent accesses to
	shared memory.

  \item[{\hyperref[sec:docs:Quick-reference guide to the Linux-kernel memory model]{
	  \path{cheatsheet.txt}}}]
	Quick-reference guide to the Linux-kernel memory model.

  \item[{\hyperref[sec:docs:Control dependencies]{
	  \path{control-dependencies.txt}}}]
	Guide to preventing compiler optimizations from destroying
	your control dependencies.

  \item[{\hyperref[sec:docs:Explanation of the Linux-Kernel Memory Consistency Model]{
	  \path{explanation.txt}}}]
	Detailed description of the memory model.

  \item[{\hyperref[sec:docs:Glossaries to LKMM]{\path{glossary.txt}}}]
	Brief definitions of LKMM-related terms.

  \item[{\hyperref[sec:docs:Herd representation]{
	  \path{herd-representation.txt}}}]
	The (abstract) representation of the Linux-kernel concurrency
	primitives in terms of events.

  \item[{\hyperref[sec:docs:Linux-kernel memory model litmus tests]{
	  \path{litmus-tests.txt}}}]
	The format, features, capabilities, and limitations of the litmus
	tests that LKMM can evaluate.

  \item[{\hyperref[sec:docs:Locking]{\path{locking.txt}}}]
	Rules for accessing lock-protected shared variables outside of
	their corresponding critical sections.

  \item[{\hyperref[sec:docs:Overview of memory-ordering operations]{
	  \path{ordering.txt}}}]
	Overview of the Linux kernel's low-level memory-ordering
	primitives by category.

  \item[{\hyperref[sec:docs:Recipes]{\path{recipes.txt}}}]
	Common memory-ordering patterns.

  \item[{\hyperref[sec:docs:Quick-reference guide to the Linux-kernel memory model]{
	  \path{references.txt}}}]
	Background information.

  \item[{\hyperref[sec:docs:Quick and simple guide to LKMM]{
	  \path{simple.txt}}}]
	Starting point for someone new to Linux-kernel concurrency.
	And also a reminder of the simpler approaches to concurrency!
\end{description}
